%% kapa: exam=1 round=31 year=2020 period=2 form=A subject=accounting subject_ko=회계학 seq=21 number=61
%% source: Q-Net 감정평가사 기출문제 (cst00309) artlSeq=5207219 file=extracted/1008/5207219/A형/A형/2020년도 제31회 감정평가사 1차 시험_2교시_A형.pdf
%% conversion: kordoc 4.24.51 -> tools/gen1.py (auto); review=proofread-sonnet; flags=none
\begin{question}{61}
(주)감평의 20x1년도 현금흐름표상 영업에서 창출된 현금(영업으로부터 창출된 현금)은 ￦100,000이다. 다음 자료를 이용하여 계산한 (주)감평의 20x1년 법인세비용차감전순이익 및 영업활동순현금흐름은? (단, 이자지급 및 법인세 납부는 영업활동으로 분류한다.)
\begin{framedtext}
\begin{tabular*}{\linewidth}{@{}l@{\extracolsep{\fill}}r@{\hspace{2em}}l@{\extracolsep{\fill}}r@{}}
매출채권손상차손 & ￦500 & 매출채권(순액) 증가 & ￦4,800 \\
감가상각비 & 1,500 & 재고자산(순액) 감소 & 2,500 \\
이자비용 & 2,700 & 매입채무 증가 & 3,500 \\
사채상환이익 & 700 & 미지급이자 증가 & 1,000 \\
법인세비용 & 4,000 & 미지급법인세 감소 & 2,000
\end{tabular*}
\end{framedtext}
{\par\noindent\hspace*{2em}\makebox[4cm][c]{\uline{법인세비용차감전순이익}}\makebox[4cm][c]{\uline{영업활동순현금흐름}}\par}
\begin{choices}
  \item \makebox[4cm][c]{￦94,800}\makebox[4cm][c]{￦92,300}
  \item \makebox[4cm][c]{￦95,300}\makebox[4cm][c]{￦92,300}
  \item \makebox[4cm][c]{￦96,800}\makebox[4cm][c]{￦95,700}
  \item \makebox[4cm][c]{￦97,300}\makebox[4cm][c]{￦95,700}
  \item \makebox[4cm][c]{￦98,000}\makebox[4cm][c]{￦107,700}
\end{choices}
\end{question}
\answer{2}
